% ── Chapter 2: Instruction ────────────────────────────────────
% The imposed lattice. Arguments as clockwork.
% Ordered on the surface; subtle fractures underneath.

\chapter{Instruction}

\strandmarker{Guyon}{France, 1655}

\lettrine{T}{he} tutor arrived on Tuesdays.

He was a precise man, in the way that people become precise when
precision has been the condition of their survival. He had
survived examinations, disputations, the slow institutional
weight of a formation designed to produce men who thought in
exactly the right ways. What it had produced, in his case, was
someone who could move through an argument without losing the
thread, and who had mistaken this ability for understanding.

He was not unkind. That was important, and she acknowledged it.
He explained carefully. He repeated what she had not grasped on
the first attempt, adjusting the angle of approach but not the
destination. He had a small supply of analogies, which he
deployed in a fixed rotation, and which returned, with minor
variations, every three or four weeks.

She learned quickly.

This was the difficulty.

When she learned quickly, he moved forward, and the structure
elaborated, and she moved forward with it---but she was not
always certain that what she had done was understand, as opposed
to memorize the appearance of understanding. The steps were
clear. She could reproduce them. She could explain, if asked, why
each one followed from the one before.

Whether the whole thing added up to something she could not say.

\scenebreak

There was a procedure for questions.

Questions were permitted, encouraged even, provided they emerged
naturally from the material and pointed toward the next step in
the sequence. This kind of question had an answer already
contained within the system, waiting to be drawn out. Asking it
was not really asking---it was a form of movement, a way of
advancing.

The other kind of question was not forbidden. It simply had
nowhere to go.

She had tried one, once, in an early lesson. She had asked not
about the argument but about what the argument was for---what it
was supposed to produce, in the person who completed it. Whether
the conclusion was meant to be believed, or only acknowledged.

The tutor had paused.

Then he had reformulated the question into something he could
answer, and answered it, and moved on.

She had watched this happen and understood, with a clarity she
could not yet express, that the system was not designed to
accommodate what she had asked. Not because the question was
wrong. Because it was addressed to something outside the
permitted space.

\scenebreak

By the time she was twelve she could follow any argument she was
given. She could identify where a proof turned, where a
distinction was introduced, where a conclusion was reached and
the machinery stopped.

She could also feel, without being able to name it, the places
where the machinery stopped but something had not.

A residue. Not error---she had checked, and the logic held---but
a remainder that the logic had not consumed. A sense that the
sequence had terminated correctly while something else continued,
unaddressed, just outside the frame of what had been proven.

She did not mention this to the tutor.

She filed the sensation away, as one files things that have no
category yet, against the possibility that a category might
eventually present itself.

It would be years before she understood that the residue was not
hers alone.

That the structure produced it necessarily.

That every argument that closed left something open.
